% =====================================================================
\section{Machine Learning Techniques for Candidate Screening}
\label{sec:ml}
% =====================================================================

\subsection{Taxonomy of ML Paradigms in Screening}

Machine learning enters an evidence-based screening system in exactly one place:
measuring properties of the evidence a session has accumulated. It does not
decide, and it is not trained on hiring outcomes. That restriction is what makes
the remainder of this section a measurement problem rather than a prediction
problem about people. Table~\ref{tab:ml-paradigms} states the paradigms and where
each is admissible.

\begin{table*}[t]
\caption{ML paradigms in candidate screening and their admissibility under an
evidence-audit design.}
\label{tab:ml-paradigms}
\small
\begin{tabular}{@{}p{2.4cm}p{5.2cm}p{3.2cm}p{4.6cm}@{}}
\toprule
Paradigm & Typical screening use & Admissible here? & Reason \\
\midrule
Supervised, outcome-labelled & Predict hire / no-hire from historical decisions &
\textbf{No} & Inherits whatever the historical decisions encoded; this is the
failure the 2018 Amazon case made famous \\
Supervised, evidence-labelled & Predict whether a session gathered enough
material to judge a claim & Yes & Target is a property of the interview, not of
the person \\
Relational / matching & Score a résumé against a posting & With care & Pooled
metrics conflate matching with candidate ranking (\S\ref{sec:withinquery}) \\
Unsupervised & Cluster claims, deduplicate skills & Yes & Structural, no verdict
produced \\
Self-supervised & Capture-quality prediction from embedding stability & Yes & No
human annotation, no person-level label \\
Reinforcement learning & Adaptive question selection & Offline only & Online
exploration means deliberately asking a real candidate a question believed to be
poor \\
\bottomrule
\end{tabular}
\end{table*}

\subsection{Representative Models and Use Cases}

\begin{figure*}[t]
  \centering
  \includegraphics[width=\textwidth]{fig13_ml_roles.pdf}
  \caption{Model families evaluated in this work and the role each plays in the
  comparison.}
  \label{fig:mlroles}
\end{figure*}

Figure~\ref{fig:mlroles} shows the families compared, and
Figure~\ref{fig:mlevolution} places them against the evolution of screening
automation. Table~\ref{tab:ml-studies} summarises how the literature has applied
these families.

\begin{figure*}[t]
  \centering
  \includegraphics[width=\textwidth]{fig12_ml_evolution.pdf}
  \caption{Evolution of screening automation, from lexical matching to grounded
  evidence audit.}
  \label{fig:mlevolution}
\end{figure*}

\begin{table*}[t]
\caption{Methodological studies bearing on model comparison in screening-like
tasks, and the role each plays in this article's protocol.}
\label{tab:ml-studies}
\small
\begin{tabular}{@{}p{1.5cm}p{4.2cm}p{4.4cm}p{5.0cm}@{}}
\toprule
Refs & Technique / contribution & Focus & Role in this work \\
\midrule
\cite{demsar2006statistical} & Friedman test with post-hoc procedures &
Comparing many classifiers & Adopted as the significance protocol \\
\cite{nadeau2003inference} & Corrected variance for overlapping folds &
Generalisation-error inference & Motivates non-parametric paired tests \\
\cite{dacrema2019progress} & Reproduction under controlled protocol & Neural
recommenders vs.\ heuristics & Template for a null result that is informative \\
\cite{musgrave2020reality} & Evaluation-practice audit & Deep metric learning &
Precedent for attributing gains to protocol \\
\cite{gorman2019splits} & Randomly regenerated splits & Standard-split fragility
& Motivates ten grouped folds over one split \\
\cite{kapoor2023leakage} & Eight-type leakage taxonomy & ML-based science &
One category quantified in \S\ref{sec:splitscheme} \\
\cite{grinsztajn2022tabular} & Trees vs.\ deep nets, 45 datasets & Tabular
benchmarks & Context for family-level spread \\
\cite{pedregosa2011scikit} & Reference implementations & Toolkit & All nine
models are its estimators \\
\bottomrule
\end{tabular}
\vspace{2pt}
\footnotesize
Datasets, baselines, and evaluation metrics differ across the cited papers; the
table records each work's role, not a performance comparison.
\end{table*}

\subsection{A Nine-Model, Six-Family Comparison}
\label{sec:ninemodel}

We compare nine classifiers spanning six model families on a public
résumé--posting fit benchmark~\cite{resumejobfit2024} of 7{,}910 balanced pairs
drawn from 643 résumés and 351 postings. Every experiment uses ten-fold
\texttt{GroupKFold} grouped by job posting, so each fold's test postings are
absent from its training data. Two
quantities are fitted inside each fold on training rows only --- the IDF corpus
statistics and the feature scaler --- because fitting either over the pooled
corpus would leak the test fold's term distribution into training. Identical fold
assignments are reused across every model, so all comparisons are paired.

Models run at fixed modest hyper-parameters with no per-model tuning budget;
\S\ref{sec:tuning} tests whether that choice matters. Scale-sensitive models
receive standardised features and tree ensembles receive raw features, since
standardising a scale-invariant model would compare preprocessing rather than
algorithms.

\begin{table*}[t]
\caption{Nine classifiers on résumé--posting fit, ten-fold \texttt{GroupKFold} by
posting. Efficiency columns are per-fold means.}
\label{tab:ninemodel}
\small
\begin{tabular}{@{}lcccrrrr@{}}
\toprule
Model & AUROC & Accuracy & Macro F1 & ECE & Fit (s) & Latency ($\mu$s/row) & Size (KB) \\
\midrule
MLP                      & \textbf{0.696 $\pm$ 0.031} & 0.634 $\pm$ 0.021 & 0.658 $\pm$ 0.038 & 0.057 & 0.44 & 0.5 & 74.7 \\
Gradient boosting        & 0.689 $\pm$ 0.028 & \textbf{0.636 $\pm$ 0.024} & 0.655 $\pm$ 0.041 & 0.047 & 2.44 & 1.9 & 393.3 \\
Random forest            & 0.688 $\pm$ 0.030 & 0.628 $\pm$ 0.026 & 0.647 $\pm$ 0.042 & 0.047 & 0.59 & 88.2 & 24{,}732.6 \\
SVM (RBF)                & 0.686 $\pm$ 0.028 & 0.635 $\pm$ 0.023 & 0.651 $\pm$ 0.039 & 0.053 & 4.64 & 195.0 & 526.8 \\
Extra trees              & 0.685 $\pm$ 0.030 & 0.632 $\pm$ 0.026 & 0.641 $\pm$ 0.044 & \textbf{0.046} & 0.39 & 84.9 & 63{,}415.9 \\
$k$-NN                   & 0.680 $\pm$ 0.031 & 0.628 $\pm$ 0.029 & 0.639 $\pm$ 0.041 & 0.057 & 0.01 & 35.5 & 1{,}272.8 \\
Hist.\ gradient boosting & 0.679 $\pm$ 0.029 & 0.626 $\pm$ 0.020 & 0.643 $\pm$ 0.037 & 0.057 & 0.41 & 7.3 & 591.8 \\
Gaussian NB              & 0.667 $\pm$ 0.033 & 0.620 $\pm$ 0.027 & 0.633 $\pm$ 0.051 & 0.117 & 0.00 & 0.3 & 1.1 \\
Logistic regression      & 0.665 $\pm$ 0.022 & 0.616 $\pm$ 0.020 & 0.637 $\pm$ 0.042 & 0.052 & 0.01 & 0.2 & \textbf{0.9} \\
\midrule
\emph{Fixed-threshold baseline} & --- & \emph{0.513 $\pm$ 0.028} & --- & --- & --- & --- & --- \\
\bottomrule
\end{tabular}
\end{table*}

The nine span 0.031 AUROC from best to worst, and every model's fold-to-fold
standard deviation is of the same order as that entire spread. The Friedman test
rejects the hypothesis that all nine perform identically ($\chi^2 = 38.35$,
$p = 6.5 \times 10^{-6}$), so the models are not interchangeable --- but the
post-hoc picture is what governs a ranking claim. Comparing the leader against
each other model with a paired Wilcoxon signed-rank test over the ten folds and
applying Holm--Bonferroni correction, the leading model is \textbf{statistically
indistinguishable from all four tree ensembles}, separating only from the linear,
probabilistic, instance-based, and kernel baselines, where the largest effect is
three AUROC points.

A study reporting only the ranked table would license the claim ``the MLP is the
best model for résumé--job fit.'' The significance analysis does not support it.
Note also the non-monotonicity across metrics: the MLP leads on AUROC, gradient
boosting on accuracy, extra trees on calibration. Which model ``wins'' depends on
which column is placed first.

\paragraph{Establishing that the protocol has power.}
A comparison that finds no difference is uninformative unless the protocol can be
shown to detect a difference when one exists. We therefore constructed a
synthetic control with a pre-specified generative process --- twelve features,
eight carrying linear weights, one interaction term, one quadratic term, four
weightless, and a 7\% label-flip rate imposing a Bayes-optimal accuracy of 0.925
--- and ran it through the identical protocol. There the same nine models
separate decisively ($\chi^2 = 55.42$, $p = 3.7 \times 10^{-9}$), with the best
recovering 94.9\% of achievable accuracy.

\begin{figure*}[t]
  \centering
  \includegraphics[width=\textwidth]{fig1_convergence_and_power.pdf}
  \caption{Per-fold distributions over ten grouped folds. \textbf{(a)} The nine
  models on the real screening task; the boxes overlap almost completely.
  \textbf{(b)} The same nine models, same folds, same code path, on the synthetic
  control, against its Bayes-optimal ceiling. Models appear in the same order in
  both panels. The contrast is the power analysis: an instrument reporting no
  difference in (a) is shown in (b) to be capable of reporting one.}
  \label{fig:power}
\end{figure*}

Figure~\ref{fig:power} places the two side by side.
\emph{The synthetic task is not a résumé task and its accuracy figures are not
screening performance.} It exists to characterise the protocol's power, and every
artifact it emits is flagged as synthetic. Its role is that of a positive control
in an assay: an instrument that reports nothing must be shown capable of
reporting something. The comparison licenses one inference --- when learnable
signal is present the protocol resolves it clearly, and on the résumé task it
does not --- so the null is a property of the task, not a limitation of the
study.

\subsection{Separating Matching from Candidate Quality}
\label{sec:withinquery}

A pooled metric rewards two distinct capabilities and cannot separate them. A
model that knows which \emph{résumés} tend to be labelled fit will score well,
and so will a model that knows which \emph{postings} a given résumé fits. Only
the second is matching.

Varying the representation while holding the protocol fixed localises where
performance comes from. A probe given only the résumé embedding --- which cannot
be matching, since it never sees the thing being matched against --- reaches
0.660 AUROC, 87\% of the best model's score and statistically indistinguishable
from a single cosine-similarity feature. Its mirror, given only the posting,
reaches 0.546, barely above chance. The best representation overall is plain
concatenation of the two embeddings at 0.755, which is to say that the
representation most enabling résumé recognition scores highest.

\begin{figure}[t]
  \centering
  \includegraphics[width=\columnwidth]{fig2_representation_ablation.pdf}
  \caption{Best AUROC per representation, identical folds throughout. The
  highlighted bars are one-sided probes that cannot be performing matching: one
  never sees the posting, the other never sees the résumé.}
  \label{fig:ablation}
\end{figure}

Figure~\ref{fig:ablation} shows the ordering. That evidence pointed us toward a
strong conclusion --- that the benchmark's labels were not relational at all ---
and a direct test refuted it. Decomposing
label variance shows \textbf{80.6\% is within-résumé}: the same résumé labelled
fit for one posting and not for another. To measure whether models access that
variance we restrict evaluation to pairs of rows sharing a résumé, one positive
and one negative, and ask whether the model ranks the positive posting higher.
Because both items share a résumé, every résumé-level property cancels exactly,
and chance is 0.5 by construction.

\begin{table}[t]
\caption{Pairwise ranking accuracy. Random forest, identical folds, chance
$= 0.5$, 4{,}112 within-résumé pairs.}
\label{tab:withinquery}
\small
\begin{tabular}{@{}lccc@{}}
\toprule
Representation & Between & Within & Binomial $p$ \\
\midrule
Relational (10-d)    & 0.6905 & \textbf{0.6333} & $3.8 \times 10^{-66}$ \\
Interaction (1536-d) & \textbf{0.7372} & 0.6240 & $2.0 \times 10^{-57}$ \\
\bottomrule
\end{tabular}
\end{table}

\begin{figure}[t]
  \centering
  \includegraphics[width=\columnwidth]{fig03_matching_isolation.pdf}
  \caption{Pairwise ranking accuracy, chance $= 0.5$. Moving to the richer
  representation raises between-résumé ranking by 0.047 while leaving
  within-résumé matching unchanged ($-0.009$, $p = 0.445$).}
  \label{fig:withinquery}
\end{figure}

\textbf{Models do match.} Both representations rank within-résumé pairs well
above chance by every test applied. A model given one résumé and two postings
picks the fitting one roughly 63\% of the time.

\textbf{But the pooled gain does not come from matching.} Moving from the ten
engineered features to the 1536-dimensional interaction encoding raises
between-résumé ranking by $+0.047$ and pooled AUROC by $+0.041$, while
within-résumé matching changes by $-0.009$, which a paired Wilcoxon test over the
ten folds does not distinguish from zero ($p = 0.445$). We state this as
\emph{the gain does not transfer}, not as \emph{matching degrades}: the data
support the absence of an improvement, not the presence of a loss.

This is the section's central measurement. The representation change a
leaderboard would record as a clear improvement bought a better ordering of
résumés against each other and left the model's ability to tell which posting
suits a candidate where it was. Nothing in the pooled metric reveals which of the
two capabilities moved.

\paragraph{The diagnostic generalises.}
Nothing in the construction is specific to résumés. Any benchmark whose instances
are (entity, item) pairs with recurring entities --- candidate--job,
patient--treatment, user--item, student--question, query--document --- admits the
same test: condition on the entity, rank within it, compare against the pooled
number. It needs no extra labels and one pass over predictions already computed.
Where the two diverge, the pooled number is measuring something other than the
advertised capability. For algorithmic hiring specifically, this is the
measurement that distinguishes a system that ranks matches from one that ranks
people --- precisely the target substitution Raghavan et
al.~\cite{raghavan2020hiring} warn carries the ethical and legal weight.

\subsection{Robustness of the Comparison}
\label{sec:splitscheme}

\begin{table}[t]
\caption{Split-scheme comparison. Identical features and model throughout.}
\label{tab:splits}
\small
\begin{tabular}{@{}llc@{}}
\toprule
Scheme & Held out from training & AUROC \\
\midrule
S1 $k$-fold, no grouping & nothing & \textbf{0.7137 $\pm$ 0.016} \\
S3 GroupKFold by résumé  & every test résumé & 0.7032 $\pm$ 0.031 \\
S2 GroupKFold by posting & every test posting & 0.6894 $\pm$ 0.016 \\
\bottomrule
\end{tabular}
\end{table}

Two objections bear on the interpretation, and both are testable.

\emph{Does the protocol still permit résumé leakage?}
Table~\ref{tab:splits} varies only the splitting scheme. The ungrouped scheme a
conventional bake-off would use scores 0.024 AUROC above the posting-grouped
protocol we adopt, so the naive protocol inflates, modestly. But closing the
résumé channel does not collapse performance --- it raises it. We expected the
résumé-grouped scheme to fall sharply if the model's apparent skill were résumé
recognition; it scores 0.014 \emph{above} the posting-grouped scheme. A
strong-memorisation account predicts the opposite sign and is refuted. The
ordering instead tracks which side of the pair is novel, and generalising to a
new posting is the harder problem --- which matters in deployment, where every
new requisition presents exactly that condition.

\emph{Did binarising a three-class label destroy the signal?}
The cleanest available contrast, Good Fit against No Fit with the ambiguous
middle class removed, does \emph{not} exceed the binarisation used
(0.6810 vs.\ 0.6894), and its fold variance triples. Binarisation is not
concealing signal.

\subsection{Does Tuning Change the Conclusion?}
\label{sec:tuning}

The claim that nine models are statistically indistinguishable is vulnerable to
the objection that they were left untuned. We tested it with a design that could
have refuted the claim: a nested cross-validation whose outer loop is the same
ten folds and whose inner three-fold loop selects among eight configurations per
model --- an identical budget for every model, since an unequal budget compares
search effort rather than algorithms.

Tuning rescues no model. The largest gain across all nine is $+0.0077$ AUROC and
three models are marginally worse tuned than untuned, which is ordinary
nested-CV selection noise. The best-to-worst spread is essentially unchanged
($0.0311 \rightarrow 0.0317$), but the significance structure \emph{contracts}:
after Holm correction the tuned leader separates from only two of eight rivals
where the untuned leader separated from four. The two models that had been
significantly behind and gained most from tuning are no longer distinguishable
from the leader. Equal-budget tuning therefore leaves the ranking substantively
unchanged and weakens rather than strengthens the case that any model is best.

\subsection{Architectural Integration of ML}

\begin{figure}[t]
  \centering
  \includegraphics[width=0.92\columnwidth]{fig14_ml_tiers.pdf}
  \caption{Where each learned component runs. The human tier is not a fallback;
  it is the only tier that produces a decision about a person.}
  \label{fig:mltiers}
\end{figure}

Figure~\ref{fig:mltiers} shows placement. Two considerations drive it. Scoring
that must survive a network outage runs on-device, which is why the evidence
sufficiency model ships as exported coefficients rather than as a service call.
Components requiring model weights too large to ship run in the service. Training
and calibration are strictly offline and gated: the export refuses to write an
artifact unless held-out AUROC, calibration error, and Brier score all clear
stated bars, so a bad run produces no file rather than a quietly bad one.

\paragraph{Efficiency is decoupled from accuracy.}
Across the nine models, inference latency spans 0.2 to 195.0 $\mu$s per row --- a
factor of roughly 1{,}000 --- and serialised size spans 0.9 KB to 63.4 MB, a
factor of about 70{,}000. Predictive quality spans 0.031 AUROC.
Figure~\ref{fig:efficiency} shows the frontier, which is nearly flat. For an
interactive deployment the rational choice is therefore the cheapest model that
is not significantly worse, which here is a kilobyte of logistic-regression
coefficients rather than a 64-megabyte forest.

\begin{figure}[t]
  \centering
  \includegraphics[width=\columnwidth]{fig4_efficiency_pareto.pdf}
  \caption{AUROC against inference latency (log scale). Three orders of magnitude
  of latency buy 0.031 AUROC, most of which is not statistically resolvable.}
  \label{fig:efficiency}
\end{figure}

\paragraph{More data would not close the gap.}
A learning curve over the same folds rises from 0.6514 AUROC at a tenth of the
training data to 0.6878 at all of it. Quadrupling the data from a quarter to full
buys 0.0204; doubling from half to full buys 0.0076. The curve is flattening well
below any useful operating point (Figure~\ref{fig:curve}), so the limitation is
not sample size and collecting more rows of this kind would not resolve it.

\begin{figure}[t]
  \centering
  \includegraphics[width=\columnwidth]{fig5_learning_curve.pdf}
  \caption{Learning curve over the same ten folds. The curve flattens well below
  any useful operating point.}
  \label{fig:curve}
\end{figure}

\subsection{Open Challenges and Research Directions}

\begin{table}[t]
\caption{Challenges in ML for evidence-based screening.}
\label{tab:ml-challenges}
\small
\begin{tabular}{@{}p{2.0cm}p{5.4cm}@{}}
\toprule
Challenge & Impact and direction \\
\midrule
No independent corpus & Of seven public résumé--fit corpora screened, the
largest shares 100\% of its documents with the corpus under study and the
independent ones pair each résumé with exactly one posting, making the
within-query statistic undefined \\
Undocumented provenance & The benchmark carries no description, licence,
citation, or labelling procedure, so what a classifier trained on it has learned
cannot be stated \\
Evidence labels are costly & Annotating evidence sufficiency requires session
replay and multiple annotators; agreement is the model's performance ceiling and
must be reported \\
Representation ceiling & Changing representation moves AUROC about twice as far
as changing algorithm; effort is better spent there \\
\bottomrule
\end{tabular}
\end{table}

The first row is a finding rather than a limitation of effort. A benchmark whose
most-downloaded derivatives regenerate its labels by model and re-pair its
documents combinatorially is a benchmark whose apparent diversity is illusory:
work reporting evaluation across ``several résumé-fit datasets'' may be
evaluating repeatedly on one. We note also that an exact-hash contamination check
passed the contaminated corpus as clean, and only normalised matching exposed it.
Contamination audits should normalise before hashing and compare entity sets
rather than only pair sets, since re-pairing the same documents preserves
document overlap while destroying pair overlap.

% =====================================================================
\section{Deep Learning for Identity and Extraction}
\label{sec:dl}
% =====================================================================

\subsection{DL Architectures}

Two deep architectures carry load in this system, and neither is trained here.
Face detection uses SCRFD~\cite{guo2022scrfd}, whose contribution is
redistributing computation between backbone, neck, and head to detect small faces
efficiently. Face embedding uses ArcFace~\cite{deng2019arcface}, whose additive
angular margin loss produces highly discriminative 512-dimensional embeddings.
Both are consumed as pretrained components.

\begin{table*}[t]
\caption{Deep architectures in the pipeline: advantages, limitations, and use.}
\label{tab:dl-arch}
\small
\begin{tabular}{@{}p{2.4cm}p{4.4cm}p{4.4cm}p{4.2cm}@{}}
\toprule
Architecture & Advantage & Limitation & Use here \\
\midrule
SCRFD detector & Efficient across compute regimes; strong on small faces &
Detection only; no identity & Locating faces in webcam frames \\
ArcFace embedding & Angular margin yields separable identity geometry &
Threshold left to the integrator & 512-d embedding for identity continuity \\
Sentence encoder & Dense semantic comparison of free text & Similarity is not
entailment & Claim clustering and relational features \\
Language model (gated) & Fluent span proposal over messy documents & Will author
text if unconstrained & Proposes spans; the gate decides \\
\bottomrule
\end{tabular}
\end{table*}

\paragraph{Why the encoder is frozen.}
Retraining face recognition on a project-scale dataset would be strictly worse
than the pretrained model and would introduce demographic bias the project cannot
audit. The research question is not the encoder; it is what is done with the
embedding stream downstream. Relatedly, the vendor model pack ships an age and
gender classifier, and the loader specifies an allow-list that never loads it ---
a single argument a reviewer can verify, rather than a policy statement.

\subsection{Threshold Calibration and Abstention}

\begin{figure*}[t]
  \centering
  \includegraphics[width=\textwidth]{fig15_dl_identity.pdf}
  \caption{Identity pipeline. The result type is three-valued;
  \emph{Unchecked} carries a reason and no similarity value at all.}
  \label{fig:identity}
\end{figure*}

The system originally compared embeddings against an authored similarity
threshold of 0.50 --- a value arrived at by reasoning, documented honestly, and
flagged in its own source comment as uncalibrated. We replaced it with a measured
one. Using the Labeled Faces in the Wild benchmark~\cite{huang2007lfw}, an
equal-error-rate sweep on 2{,}200 training pairs and evaluation on 1{,}000
held-out pairs yields the results in Table~\ref{tab:calibration}. No weights are
trained; a single scalar is fitted.

\begin{table}[t]
\caption{Face-verification threshold before and after calibration on LFW.}
\label{tab:calibration}
\small
\begin{tabular}{@{}lcccc@{}}
\toprule
Threshold & FAR & FRR & Accuracy & AUC \\
\midrule
\textbf{0.1266} (calibrated) & 0.030 & \textbf{0.034} & \textbf{0.968} & 0.9785 \\
0.50 (authored) & 0.000 & \textbf{0.414} & 0.793 & --- \\
\bottomrule
\end{tabular}
\end{table}

\textbf{This is the article's most consequential single finding.} The prior
threshold was not arbitrary; it was reasoned, documented, and plausible. It was
also wrong in a way nobody could have known without measuring: it would have
falsely rejected \textbf{41.4\%} of genuine identity matches. In deployment,
roughly two in five legitimate re-verification checks would have failed --- and
because the system correctly refuses to treat an unmeasured check as a pass,
those candidates would have accumulated integrity flags for being themselves.
Calibration reduced false rejection by a factor of twelve at a false-acceptance
cost of 3.0\%. The lesson generalises beyond this component: in
security-adjacent systems, \emph{reasoned} and \emph{correct} are different
properties, and only one of them can be demonstrated.

\paragraph{Abstention as a hard gate.}
Figure~\ref{fig:identity} shows the three-valued result type.
\texttt{Unchecked} carries a reason and has \emph{no similarity field at all}. It
is not similarity 0.0, not null, not a default --- the value does not exist, so
no downstream comparison can read a failed measurement as a low-but-real one and
no serialisation can emit a number that was never measured. Selective prediction
of this kind is a hard constraint rather than a soft preference, in the spirit of
distribution-free abstention~\cite{vovk2005algorithmic}; a model that
\emph{usually} respects a safety constraint is a model that sometimes does not.

\paragraph{Stated limitation.}
LFW comprises public-figure photographs, not this system's candidate population.
The calibrated value is a defensible general threshold, not one validated on real
candidates, and population-specific recalibration requires consented candidate
data we have not collected.

\subsection{Deployment Challenges}

\begin{table}[t]
\caption{Challenges, impacts, and mitigations for DL components.}
\label{tab:dl-challenges}
\small
\begin{tabular}{@{}p{1.9cm}p{2.6cm}p{2.6cm}@{}}
\toprule
Challenge & Impact & Mitigation \\
\midrule
Population shift & Benchmark-calibrated thresholds may not transfer & Report the
mismatch; recalibrate under consent \\
Capture quality & Poor lighting widens within-session variance & Abstain rather
than mismatch \\
Silent model failure & A failed load can be masked by a fallback & Surface
component health in the UI; never substitute a heuristic \\
Embedding privacy & Face embeddings are invertible to recognisable faces &
Retain on-device, encrypted, with a deletion window \\
Demographic inference & Bundled classifiers infer protected attributes &
Allow-list modules; never load them \\
\bottomrule
\end{tabular}
\end{table}

% =====================================================================
\section{Synergistic Integration of Grounding, ML, and DL}
\label{sec:integration}
% =====================================================================

\subsection{Motivation for Integration}

Each family alone leaves a gap. Grounding guarantees that claims are the
candidate's own words but says nothing about whether a session gathered enough
evidence to judge them. ML measures evidence sufficiency but, without grounding,
would measure the sufficiency of text a model may have invented. DL supplies
identity continuity but produces no claim and no verdict. The combination closes
a loop that none closes alone, and Figure~\ref{fig:loop} shows its shape --- with
the human reviewer at the centre, because the loop deliberately terminates in a
person rather than in an aggregate.

\begin{figure}[t]
  \centering
  \includegraphics[width=0.92\columnwidth]{fig16_feedback_loop.pdf}
  \caption{The integrated loop. Every arrow carries typed, cited evidence; the
  reviewer is the only component that produces a decision about a person.}
  \label{fig:loop}
\end{figure}

\subsection{Evaluated Benefits and Impact}

\begin{table*}[t]
\caption{Key metrics before and after integration. Every row is a measured value
from this work, not an estimate.}
\label{tab:before-after}
\small
\begin{tabular}{@{}p{3.6cm}p{4.4cm}p{4.4cm}p{3.4cm}@{}}
\toprule
Metric & Before & After & Basis \\
\midrule
Identity false-rejection rate & 0.414 (authored threshold 0.50) &
\textbf{0.034} (calibrated 0.1266) & LFW, 1{,}000 held-out pairs \\
Paraphrase leakage & Unmeasured & \textbf{0 of 48{,}149} & Offline gate harness \\
Negation-trap leakage & Unmeasured & \textbf{0 of 51{,}761} & Offline gate harness \\
Verbatim admission & 99.18\% (425 wrongly rejected) & \textbf{100.00\%} &
Clause-splitter fix \\
Model-selection basis & Ranked table & Friedman + Holm-corrected pairwise &
Ten grouped folds \\
Matching vs.\ ranking & Conflated in one metric & Separated, 0.633 within-résumé
& 4{,}112 within-résumé pairs \\
Composite score & Present in comparable systems & \textbf{Absent by construction}
& Vocabulary ban fails the build \\
\bottomrule
\end{tabular}
\end{table*}

Table~\ref{tab:before-after} collects the measured effect of integration. Three
observations follow. First, the largest single improvement came not from a model
but from measuring a constant. Second, two rows read ``unmeasured'' before ---
the property was asserted rather than tested, which is the condition this article
argues is most common and least noticed. Third, the last row is a capability
\emph{removed}: the system cannot emit a composite score, and that absence is
enforced by a test.

\subsection{Architecture as Guarantee}

The recurring engineering move across all three families is converting a
\emph{policy} into a \emph{structure}.

\begin{table}[t]
\caption{Policies converted into structural guarantees.}
\label{tab:policy-structure}
\small
\begin{tabular}{@{}p{2.4cm}p{4.9cm}@{}}
\toprule
Policy & Structure \\
\midrule
AI must not author claims & Import boundary plus a CI architecture test \\
Unmeasured is not passed & A result type with no similarity field \\
No composite score & Vocabulary ban that fails the build \\
No demographic inference & Classifier module never loaded \\
No outcome labels & No such dataset and no code path to build one \\
Plans cannot be back-edited & Question plans are never persisted \\
Sessions must be reconstructible & Append-only event table with a sequencing
trigger \\
\bottomrule
\end{tabular}
\end{table}

Each row converts a claim that requires trust into one that can be checked in
seconds by someone who does not trust us. That is the property we would most like
other screening systems to adopt, independent of any of the specific technologies
in this article.

\subsection{Integration Challenges and Design Considerations}

\begin{table}[t]
\caption{Barriers to integration and mitigations.}
\label{tab:barriers}
\small
\begin{tabular}{@{}p{1.9cm}p{5.4cm}@{}}
\toprule
Barrier & Mitigation \\
\midrule
Interoperability & Typed contracts between surfaces; a failed measurement has its
own representation rather than a sentinel value \\
Lifecycle & Export gates refuse to write an artifact that misses its quality
bars; artifacts declare their own provenance and the loader rejects one that
misrepresents it \\
Latency budget & On-device scoring for anything on the interactive path; the
cheapest model that is not significantly worse \\
Compliance & Evidence retained per claim; no aggregate to explain because none is
produced \\
Drift & Thresholds and calibration are dated and re-measurable; nothing is
authored once and trusted indefinitely \\
\bottomrule
\end{tabular}
\end{table}

\subsection{Standards and Regulatory Compliance}

\begin{table*}[t]
\caption{Mapping regulatory obligations onto architectural mechanisms.}
\label{tab:compliance}
\small
\begin{tabular}{@{}p{3.4cm}p{5.4cm}p{6.4cm}@{}}
\toprule
Obligation & Instrument & Mechanism in this system \\
\midrule
Bias audit of the tool & NYC Local Law 144 & No composite score to audit; the
audited object is the per-claim verdict distribution \\
Candidate notification & NYC Local Law 144; Illinois AIVIA & Consent captured
per session with a recorded consent version \\
Human oversight & EU AI Act, Annex III high-risk & Report generation is a
deterministic transform containing no code that could emit a recommendation \\
Explanation and contestability & EU AI Act; HB 3773 & Every verdict carries its
evidence and character offsets into the candidate's own words \\
Non-discrimination & HB 3773; Illinois Human Rights Act & No demographic
inference, no outcome labels, no historical decision pattern to inherit \\
Traceability & All & Append-only event store; sessions reconstructible turn by
turn \\
\bottomrule
\end{tabular}
\end{table*}

The mapping in Table~\ref{tab:compliance} is possible only because the system
declines to produce an aggregate. A bias audit of a composite score must explain
the score; a bias audit of a claim audit examines the distribution of verdicts
and the evidence attached to each, which is a smaller and better-posed question.
